\documentclass[11pt,a4paper]{article}

\usepackage{fontspec}
\setmainfont{PTF55F.ttf}[BoldFont=PTF75F.ttf,ItalicFont=PTF56F.ttf,BoldItalicFont=PTF76F.ttf]
\setsansfont{PTS55F.ttf}[BoldFont=PTS75F.ttf,ItalicFont=PTS56F.ttf,BoldItalicFont=PTS76F.ttf]
\setmonofont{PTM55F.ttf}[Scale=0.92]
\usepackage[russian]{babel}
\AtBeginDocument{\shorthandoff{"}}

\usepackage[margin=2.2cm]{geometry}
\usepackage{xcolor}
\usepackage{minted}
\usepackage[most,minted]{tcolorbox}
\usepackage{booktabs}
\usepackage{tabularx}
\usepackage{enumitem}
\usepackage{titlesec}
\usepackage[colorlinks=true,linkcolor=blue!50!black,urlcolor=blue!50!black]{hyperref}

\setlength{\parindent}{0pt}
\setlength{\emergencystretch}{3em}
\setlength{\parskip}{0.55em}
\setlist{itemsep=0.15em,topsep=0.3em}
\titleformat{\section}{\Large\sffamily\bfseries}{\thesection}{0.6em}{}
\titleformat{\subsection}{\large\sffamily\bfseries}{\thesubsection}{0.6em}{}
\titleformat{\subsubsection}{\normalsize\sffamily\bfseries}{\thesubsubsection}{0.6em}{}

\definecolor{codebg}{HTML}{F6F7F9}

% Листинги: minted (нужен pygmentize и tectonic -Z shell-escape).
\usemintedstyle{friendly}
% Пробелы в листингах выводятся настоящим глифом U+0020, чтобы отступы
% сохранялись при копировании кода из PDF.
\newcommand{\copyspace}{\XeTeXglyph\XeTeXcharglyph"0020\relax}
\tcbset{codebox/.style={listing only, breakable, enhanced, colback=codebg, frame hidden,
  arc=1pt, left=4pt, right=4pt, top=1pt, bottom=1pt, before skip=0.5em, after skip=0.6em}}
\newtcblisting{py}{codebox, listing engine=minted, minted language=python, minted options={fontsize=\small,breaklines,tabsize=4,showspaces,space=\copyspace}}
\newtcblisting{shell}{codebox, listing engine=minted, minted language=console, minted options={fontsize=\small,breaklines,tabsize=4,showspaces,space=\copyspace}}
\newtcblisting{toml}{codebox, listing engine=minted, minted language=toml, minted options={fontsize=\small,breaklines,tabsize=4,showspaces,space=\copyspace}}
\newtcblisting{plaintext}{codebox, listing engine=minted, minted language=text, minted options={fontsize=\small,breaklines,tabsize=4,showspaces,space=\copyspace}}
\newcommand{\code}[1]{\texttt{#1}}

\begin{document}

\begin{center}
  {\sffamily\small Программирование на языке Python · 2 курс}\\[0.6em]
  {\sffamily\LARGE\bfseries Лабораторная работа №1}\\[0.4em]
  {\sffamily\Large Окружение, git и синтаксис Python.\\ Задача «Миллиард строк»}
\end{center}

\tableofcontents
\newpage

%=====================================================================
\section{Общие положения}

\subsection{Цель работы}

Лабораторная работа предназначена для студентов, освоивших на первом курсе язык C++ и основы
объектно-ориентированного программирования. Цель работы~— перенос имеющихся знаний на язык
Python и изучение его принципиальных отличий от C++. По итогам работы студент должен:

\begin{itemize}
  \item уметь настраивать рабочее окружение: интерпретатор, виртуальное окружение, менеджер проектов \code{uv};
  \item уметь вести историю проекта в git и публиковать репозиторий на GitHub;
  \item владеть базовым синтаксисом языка, встроенными типами данных и работой с файлами;
  \item уметь решать задачу обработки большого текстового файла средствами стандартной библиотеки.
\end{itemize}

%=====================================================================
\section{Рабочее окружение}

\subsection{Интерпретатор}

Программа на C++ проходит путь \emph{исходный код $\to$ компилятор $\to$ машинный код $\to$ запуск}.
В основной реализации Python (CPython) отдельного этапа сборки нет: интерпретатор \code{python}
компилирует исходный код в байт-код и сразу исполняет его на собственной виртуальной машине.
Как следствие, ошибки, которые в C++ выявляются при компиляции (например, опечатка в имени
переменной), в Python проявляются только при выполнении соответствующей строки.

\subsection{Пакеты и виртуальные окружения}

Сторонние библиотеки устанавливаются менеджером пакетов \code{pip} из репозитория PyPI.
Установка пакетов в системный интерпретатор приводит к конфликтам: разным проектам могут
требоваться несовместимые версии одной библиотеки, а на macOS и Linux системный Python
используется самой операционной системой.

Для изоляции проектов применяется \textbf{виртуальное окружение} (venv)~— каталог внутри проекта,
содержащий ссылку на интерпретатор и собственный набор установленных пакетов.

\begin{shell}
$ python3 -m venv .venv            # создание окружения в каталоге .venv
$ source .venv/bin/activate        # активация (macOS / Linux)
> .venv\Scripts\activate            # активация (Windows)
(.venv) $ which python             # .venv/bin/python
(.venv) $ pip install requests     # пакет устанавливается только в .venv
(.venv) $ pip freeze > requirements.txt
(.venv) $ deactivate
\end{shell}

Такой подход требует ручной активации окружения, ручного ведения списка зависимостей и отдельной
установки нужной версии интерпретатора. В курсе вместо него используется менеджер проектов \code{uv}.

\subsection{Менеджер проектов uv}\label{sec:uv}

\href{https://docs.astral.sh/uv/}{\code{uv}}~— менеджер проектов для Python, который
устанавливает интерпретаторы требуемых версий, создаёт виртуальные окружения, управляет
зависимостями и запускает код. По сути он автоматизирует действия, выполняемые вручную
с помощью \code{venv} и \code{pip}.

\subsubsection*{Установка}\phantomsection\label{sec:uv-install}
\begin{shell}
# macOS / Linux
$ curl -LsSf https://astral.sh/uv/install.sh | sh
# Windows (PowerShell)
> powershell -ExecutionPolicy ByPass -c "irm https://astral.sh/uv/install.ps1 | iex"

$ uv --version
\end{shell}

\subsubsection*{Создание проекта}
\begin{shell}
$ uv python install 3.13           # установка интерпретатора Python 3.13
$ uv init --python 3.13 lab1       # создание проекта в каталоге lab1
$ cd lab1
$ tree -a -I .git
.
├── .gitignore
├── .python-version                # версия Python для проекта
├── README.md
├── main.py
└── pyproject.toml                 # описание проекта и зависимостей
\end{shell}

Файл \code{pyproject.toml} выполняет роль, аналогичную \code{CMakeLists.txt}: в нём указываются
имя проекта, требуемая версия Python и список зависимостей.

\begin{toml}
[project]
name = "lab1"
version = "0.1.0"
requires-python = ">=3.13"
dependencies = []
\end{toml}

\subsubsection*{Запуск и зависимости}
\begin{shell}
$ uv run main.py             # создаёт .venv при необходимости и запускает скрипт
$ uv add requests            # добавление зависимости в pyproject.toml и установка
$ uv remove requests         # удаление зависимости
$ uv sync                    # синхронизация .venv с pyproject.toml и uv.lock
\end{shell}

После первого запуска в проекте появляются каталог \code{.venv} и файл \code{uv.lock}.
В lock-файле фиксируются точные версии всех зависимостей, включая транзитивные, что обеспечивает
одинаковое окружение у всех, кто работает с репозиторием.

\subsubsection*{Редактор}

Для работы подходят VS Code с расширением \emph{Python} (Microsoft) или PyCharm. В настройках
редактора в качестве интерпретатора необходимо выбрать интерпретатор из каталога \code{.venv}
проекта (в VS Code~— команда \emph{Python: Select Interpreter}); в противном случае подсказки
и запуск будут использовать другое окружение.

\subsection{Git и GitHub}\label{sec:git}

Все лабораторные работы курса сдаются через git-репозиторий; при проверке оценивается как итоговый
код, так и история его изменений. Git~— система контроля версий, которая хранит последовательность
снимков проекта (\emph{коммитов}) и позволяет просматривать изменения и возвращаться к любому
из сохранённых состояний. GitHub~— сервис для размещения git-репозиториев, доступных по ссылке.

\subsubsection*{Установка и первичная настройка}

Git устанавливается с сайта \url{https://git-scm.com} (на macOS также через \code{xcode-select --install}
или \code{brew install git}). Перед началом работы указываются имя и адрес электронной почты,
которые записываются в каждый коммит:
\begin{shell}
$ git config --global user.name "Иван Иванов"
$ git config --global user.email "ivan@example.com"
$ git config --global init.defaultBranch main
\end{shell}

\subsubsection*{Основной цикл работы}

Команда \code{uv init} создаёт git-репозиторий в каталоге проекта (при его отсутствии используется
\code{git init}). Работа с репозиторием строится по циклу: изменение файлов $\to$ просмотр
изменений $\to$ добавление изменений в индекс $\to$ создание коммита.

\begin{shell}
$ git status                       # изменённые и неотслеживаемые файлы
$ git diff                         # содержание изменений
$ git add stats.py                 # добавление файла в следующий коммит
$ git commit -m "StationStats: update и merge"
$ git log --oneline                # история коммитов
\end{shell}

Требования к истории изменений:
\begin{itemize}
  \item каждый коммит содержит одно законченное по смыслу изменение (например, <<разбор строки
        без \code{float}>>, <<чтение файла блоками>>); единственный коммит со всей работой
        не считается приемлемой историей;
  \item сообщение коммита кратко описывает сделанное изменение: \code{"StationStats: merge"},
        а не \code{"fix"} или \code{"asdf"};
  \item перед \code{git add .} проверяется вывод \code{git status}, чтобы в коммит не попали лишние файлы.
\end{itemize}

\subsubsection*{Файлы, исключаемые из репозитория}

Файл \code{.gitignore} содержит перечень путей, которые git не отслеживает; \code{uv init}
создаёт его с \code{.venv} и кешами Python.

\subsubsection*{Публикация на GitHub}

\begin{enumerate}
  \item Регистрация на \url{https://github.com}.
  \item Настройка доступа с компьютера. Наиболее простой способ~— GitHub CLI: установить \code{gh}
        (\url{https://cli.github.com}) и выполнить \code{gh auth login}. Альтернативный способ~—
        SSH-ключ: \code{ssh-keygen -t ed25519}, после чего публичный ключ
        \code{\textasciitilde/.ssh/id\_ed25519.pub} добавляется в \emph{Settings $\to$ SSH and GPG keys}.
        Пароль от учётной записи для \code{git push} не используется.
  \item Создание пустого репозитория (\emph{New repository}, без README и \code{.gitignore},
        поскольку они уже созданы локально). Репозиторий может быть публичным или приватным.
  \item Связывание локального репозитория с удалённым и отправка коммитов:
\begin{shell}
$ git remote add origin git@github.com:<логин>/python-lab1.git
$ git push -u origin main          # первая отправка; далее достаточно git push
\end{shell}
        При использовании \code{gh} шаги 3--4 выполняются одной командой:
        \code{gh repo create python-lab1 --private --source=. --push}.
\end{enumerate}

%=====================================================================
\section{Синтаксис языка}

\subsection{Структура программы}

\begin{py}
def greet(name: str) -> str:
    return f"Привет, {name}!"


if __name__ == "__main__":
    for who in ["мир", "БГУ"]:
        print(greet(who))
\end{py}

Основные отличия от C++:
\begin{itemize}
  \item \textbf{блоки кода задаются отступами}, а не фигурными скобками; двоеточие открывает блок.
        Стандартный отступ~— 4 пробела;
  \item точка с запятой в конце инструкции не ставится, типы переменных не объявляются;
  \item \code{name: str} и \code{-> str}~— \emph{аннотации типов} (type hints). Интерпретатор
        их не проверяет: они служат документацией и используются редактором для подсказок.
        Проверку соответствия типов выполняют отдельные инструменты статического анализа
        (например, \code{mypy}), которые рассматриваются в лабораторной работе №2;
  \item \code{f"...\{expr\}..."}~— f-строка, аналог \code{std::format};
  \item функция \code{main} не обязательна: код модуля выполняется сверху вниз. Условие
        \code{\_\_name\_\_ == "\_\_main\_\_"} истинно, только если файл запущен как скрипт,
        а не импортирован из другого модуля.
\end{itemize}

Стиль кода на Python описан в документе \href{https://peps.python.org/pep-0008/}{PEP~8}:
\code{snake\_case} для функций и переменных, \code{CamelCase} для классов, \code{UPPER\_CASE}
для констант.

\subsection{Переменные, объекты и ссылки}

В C++ переменная обозначает область памяти определённого типа. В Python переменная~— это
\textbf{имя, связанное с объектом}. Тип принадлежит объекту, а не имени, и одно имя может быть
повторно связано с объектом другого типа.

\begin{py}
x = 42          # имя x связано с объектом int
x = "строка"    # теперь -- с объектом str
a = [1, 2, 3]
b = a           # копирования нет: a и b -- два имени одного списка
b.append(4)
print(a)        # [1, 2, 3, 4]
c = a.copy()    # явная (поверхностная) копия
print(a == c, a is c)   # True False: == сравнивает значения, is -- идентичность
\end{py}

Семантику переменных Python удобно представлять через \code{std::shared\_ptr<Object>} в C++:
присваивание копирует ссылку, а не объект. Память освобождается автоматически (подсчёт ссылок
и сборщик циклов). Аргументы функций также передаются по ссылке.

Объекты делятся на \textbf{изменяемые} (\code{list}, \code{dict}, \code{set}, большинство
пользовательских классов) и \textbf{неизменяемые} (\code{int}, \code{float}, \code{str},
\code{tuple}, \code{bytes}, \code{frozenset}). Содержимое неизменяемого объекта изменить нельзя~—
можно только создать новый объект:

\begin{py}
s = "abc"
s += "d"        # создаётся новая строка "abcd", имя s связывается с ней
# s[0] = "X"    # TypeError: 'str' object does not support item assignment
\end{py}

\subsection{Встроенные типы}

\begin{tabularx}{\linewidth}{@{}lXl@{}}
\toprule
\textbf{Тип} & \textbf{Описание} & \textbf{Аналог в C++} \\
\midrule
\code{int}   & целое \emph{произвольной} длины, без переполнения & \code{long long} без границ \\
\code{float} & число с плавающей точкой двойной точности & \code{double} \\
\code{bool}  & \code{True} / \code{False} (подкласс \code{int}) & \code{bool} \\
\code{str}   & неизменяемая строка символов Unicode & \code{std::u32string} \\
\code{bytes} & неизменяемая последовательность байтов & \code{std::vector<uint8\_t>} \\
\code{None}  & отсутствие значения, единственный объект своего типа & \code{nullptr} / \code{std::nullopt} \\
\code{list}  & динамический массив ссылок на объекты & \code{std::vector<Object*>} \\
\code{tuple} & неизменяемая последовательность фиксированной длины & \code{std::tuple} \\
\code{dict}  & хеш-таблица, сохраняющая порядок вставки & \code{std::unordered\_map} \\
\code{set}   & хеш-множество & \code{std::unordered\_set} \\
\bottomrule
\end{tabularx}

\begin{py}
7 / 2       # 3.5   -- деление всегда даёт float
7 // 2      # 3     -- целочисленное деление с округлением вниз
-7 // 2     # -4    -- в C++ результат равен -3
-7 % 2      # 1     -- знак остатка совпадает со знаком делителя
2 ** 10     # 1024  -- возведение в степень
int("42"), float("3.14"), str(42)   # преобразования типов
len("Минск")                        # 5 символов
len("Минск".encode("utf-8"))        # 10 байт в кодировке UTF-8
\end{py}

\subsection{Строки}

\begin{py}
s = "Hello, World"
s.lower(), s.upper()            # 'hello, world', 'HELLO, WORLD'
s.split(", ")                   # ['Hello', 'World']
"key=value".partition("=")      # ('key', '=', 'value')
"  text \n".strip()             # 'text'
", ".join(["a", "b", "c"])      # 'a, b, c'
s.startswith("He"), "or" in s   # True True
s.find("o"), s.replace("l", "L")
f"{3.14159:.2f}", f"{42:>6}"    # '3.14', '    42'
\end{py}

Поскольку строки неизменяемы, операция \code{s += x} в цикле на каждой итерации создаёт новую
строку. Для объединения большого числа фрагментов их накапливают в списке и объединяют вызовом
\code{"".join(parts)}.

\subsection{Коллекции}

\begin{py}
xs = [3, 1, 2]
xs.append(5)            # добавление в конец, амортизированно O(1)
xs.pop()                # удаление с конца, O(1)
xs.insert(0, 10)        # вставка, O(n)
xs.sort()               # сортировка на месте
ys = sorted(xs, reverse=True)   # новый отсортированный список
xs[0], xs[-1]           # первый и последний элементы
xs[1:3], xs[::-1]       # срезы [start:stop:step] создают копию части списка

point = (1.0, 2.0)      # кортеж
x, y = point            # распаковка
a, b = b, a             # обмен значений без временной переменной

ages = {"Анна": 19, "Борис": 20}
ages["Вера"] = 18       # вставка или перезапись
ages.get("Глеб")        # None, если ключ отсутствует (ages["Глеб"] -- KeyError)
ages.get("Глеб", 0)     # значение по умолчанию
"Анна" in ages          # проверка наличия ключа, в среднем O(1)
for name, age in ages.items():
    print(name, age)

seen = {1, 2, 3}
seen.add(2)             # множество не изменится
seen & {2, 5}, seen | {5}   # пересечение, объединение
\end{py}

Ключом словаря и элементом множества может быть только \emph{хешируемый} объект~— как правило,
неизменяемый: \code{int}, \code{str}, \code{tuple} из хешируемых элементов. Список не может
быть ключом словаря.

\subsection{Управляющие конструкции}

\begin{py}
if x > 0 and not done:          # and, or, not вместо &&, ||, !
    ...                         # ... (или pass) -- пустое тело блока
elif x == 0:
    ...
else:
    ...

for i in range(10):             # 0..9; в общем виде range(start, stop, step)
    ...
for i, item in enumerate(items):    # индекс и элемент
    ...
for a, b in zip(xs, ys):        # параллельный обход двух последовательностей
    ...

while n > 1:
    n //= 2                     # операторов ++ и -- нет: n += 1

label = "чёт" if n % 2 == 0 else "нечет"    # условное выражение

match command.split():          # сопоставление с образцом (Python 3.10+)
    case ["go", direction]:
        move(direction)
    case ["quit" | "exit"]:
        stop()
    case _:
        print("неизвестная команда")
\end{py}

Цикл \code{for} в Python всегда перебирает элементы коллекции, аналогично range-based for
в C++11; конструкции вида \code{for (int i = 0; i < n; ++i)} нет. Если помимо элемента требуется
его индекс, используется \code{enumerate}, а не \code{range(len(xs))}.

\subsection{Генераторы коллекций}

Генераторы коллекций (comprehensions)~— сокращённая запись построения коллекций:
\begin{py}
squares = [x * x for x in range(10)]
evens   = [x for x in xs if x % 2 == 0]
lengths = {word: len(word) for word in words}       # словарь
letters = {ch for ch in "abracadabra"}              # множество
total   = sum(x * x for x in range(10))             # без промежуточного списка
\end{py}

\subsection{Функции}

\begin{py}
def power(base: float, exp: int = 2) -> float:      # значение по умолчанию
    return base ** exp

power(3), power(3, 3), power(exp=3, base=2)         # позиционные и именованные аргументы

def mean(*values: float) -> float:                  # произвольное число аргументов
    return sum(values) / len(values)

mean(1, 2, 3)

def connect(host: str, *, timeout: float = 1.0, **options) -> None:
    # аргументы после * передаются только по имени; **options -- словарь прочих
    print(host, timeout, options)

connect("bsu.by", timeout=5, retries=3)

def min_max(xs: list[int]) -> tuple[int, int]:      # возврат нескольких значений
    return min(xs), max(xs)

lo, hi = min_max([3, 1, 4])
\end{py}

Функции являются объектами: их можно присваивать переменным, передавать в другие функции
и создавать в выражениях с помощью \code{lambda}:
\begin{py}
words.sort(key=len)                         # сортировка по длине
pairs.sort(key=lambda p: (p[1], p[0]))      # по второму, затем по первому элементу
\end{py}

Значение аргумента по умолчанию вычисляется \textbf{один раз}~— при определении функции.
Изменяемое значение по умолчанию поэтому оказывается общим для всех вызовов:
\begin{py}
def add(item, bucket=[]):       # ошибка
    bucket.append(item)
    return bucket

add(1); add(2)                  # второй вызов вернёт [1, 2]

def add(item, bucket=None):     # корректный вариант
    if bucket is None:
        bucket = []
    bucket.append(item)
    return bucket
\end{py}

\subsection{Исключения}

\begin{py}
try:
    value = int(text)
except ValueError as e:
    print(f"не число: {e}")
    value = 0
else:
    print("разбор выполнен без ошибок")   # если исключения не было
finally:
    print("выполняется всегда")

if count < 0:
    raise ValueError(f"count должен быть >= 0, получено {count}")
\end{py}

Часто используемые встроенные исключения: \code{ValueError}, \code{TypeError},
\code{KeyError}, \code{IndexError}, \code{FileNotFoundError}, \code{ZeroDivisionError}.
Пользовательские исключения наследуются от \code{Exception}.

В отличие от C++, где исключения принято использовать лишь для исключительных ситуаций,
в Python они являются штатным механизмом управления потоком выполнения: например, цикл
\code{for} завершается по исключению \code{StopIteration}. Стиль <<выполнить и обработать
ошибку>> (EAFP, \emph{easier to ask for forgiveness than permission}) считается в Python нормой.

\subsection{Работа с файлами}

\begin{py}
with open("data.txt", encoding="utf-8") as f:       # текстовый режим
    for line in f:                                  # построчное чтение
        print(line.rstrip("\n"))
# здесь файл уже закрыт, даже если внутри блока возникло исключение

with open("out.txt", "w", encoding="utf-8") as f:
    f.write("строка\n")

with open("data.txt", "rb") as f:                   # бинарный режим
    chunk = f.read(1 << 20)                         # чтение до 1 МиБ -> bytes
    text = chunk.decode("utf-8")                    # bytes -> str
\end{py}

Конструкция \code{with} (контекстный менеджер) реализует идиому RAII: ресурс гарантированно
освобождается при выходе из блока. В текстовом режиме байты файла декодируются в \code{str};
кодировку следует указывать явно (\code{encoding="utf-8"}), поскольку кодировка по умолчанию
зависит от операционной системы.

\subsection{Соответствие конструкций C++ и Python}

\begin{tabularx}{\linewidth}{@{}>{\ttfamily}X>{\ttfamily}X@{}}
\toprule
\textrm{\textbf{C++}} & \textrm{\textbf{Python}} \\
\midrule
std::vector<int> v\{1, 2\}; & v = [1, 2] \\
v.push\_back(3); v.size(); & v.append(3); len(v) \\
std::unordered\_map<K, V> m; & m = \{\} \\
if (m.count(k)) & if k in m: \\
for (auto\& x : v) & for x in v: \\
for (int i = 0; i < n; ++i) & for i in range(n): \\
a \&\& b, a || b, !a & a and b, a or b, not a \\
x ? a : b & a if x else b \\
nullptr & None \\
std::cout << x << '\textbackslash n'; & print(x) \\
std::getline(std::cin, s); & s = input() \\
std::stoi(s), std::to\_string(x) & int(s), str(x) \\
// комментарий & \# комментарий \\
throw std::runtime\_error("..."); & raise RuntimeError("...") \\
try \{\} catch (const E\& e) \{\} & try: ... except E as e: ... \\
\bottomrule
\end{tabularx}

%=====================================================================
\section{Задание}

Работа состоит из двух частей. Результат сдаётся в виде git-репозитория
на GitHub (раздел~\ref{sec:git}); преподавателю направляется ссылка на репозиторий.
Тот же репозиторий используется в последующих лабораторных работах.

\subsection{Часть 1. Создание проекта}

\begin{enumerate}
  \item \hyperref[sec:uv-install]{Установить \code{uv}} (раздел~\ref{sec:uv}) и Python~3.13 (\code{uv python install 3.13}).
  \item Создать проект командой \code{uv init --python 3.13}, скопировать в него каталог \code{data/}
        из материалов работы и выполнить первый коммит.
  \item Дополнить \code{.gitignore} так, чтобы в репозиторий не попадали каталог \code{.venv}
        и сгенерированные файлы с данными.
  \item Создать репозиторий на GitHub (публичный либо приватный с доступом для преподавателя)
        и отправить в него проект. В ходе дальнейшей работы изменения фиксируются
        отдельными коммитами и регулярно отправляются командой \code{git push}.
  \item Итоговая структура проекта:
\begin{plaintext}
lab1/
├── .gitignore
├── .python-version
├── pyproject.toml
├── uv.lock
├── data/              # материалы работы (генератор, примеры, check.py)
├── README.md          # инструкция по запуску
└── main.py            # решение задачи из части 2
\end{plaintext}
  \item Все функции снабжаются аннотациями типов; имена соответствуют PEP~8.
\end{enumerate}

\subsection{Часть 2. Задача «Миллиард строк»}

В 2024 году в Java-сообществе проводился конкурс
\href{https://github.com/gunnarmorling/1brc}{One Billion Row Challenge}: требовалось как можно
быстрее обработать текстовый файл из миллиарда строк с измерениями температуры. Задача проста
по формулировке, однако при таком объёме данных каждая лишняя операция в цикле обработки
обходится в секунды. В данной работе требуется корректное решение задачи на Python.
Решение может быть любой сложности, в том числе с использованием многопоточности и различных механизмов
операционной системы, но без сторонних библиотек. Данные темы также будут затронуты только в последующих
лабораторных работах.

\subsubsection*{Входные данные}
f
Текстовый файл в кодировке UTF-8; каждая строка имеет вид \code{<станция>;<температура>}:
\begin{plaintext}
Minsk;7.3
Hamburg;12.0
Zürich;-3.4
Bulawayo;8.9
Palembang;38.8
\end{plaintext}

Гарантируется, что:
\begin{itemize}
  \item название станции~— непустая строка длиной не более 100 байт в UTF-8, не содержащая
        символов \code{;} и \code{\textbackslash n};
  \item температура~— число от $-99.9$ до $99.9$ ровно с одним знаком после точки;
  \item число различных станций не превышает 10\,000;
  \item каждая строка, включая последнюю, завершается символом \code{\textbackslash n}.
\end{itemize}

\subsubsection*{Выходные данные}

Для каждой станции вычисляются минимальная, средняя и максимальная температура. Результат
выводится по одной строке на станцию в формате \code{<станция>;<min>;<mean>;<max>}:
\begin{plaintext}
Abha;-0.8;18.3;48.0
Abidjan;-5.2;24.4;54.4
...
Zürich;-22.6;9.8;36.6
İzmir;-7.3;17.2;44.7
\end{plaintext}
\begin{itemize}
  \item станции упорядочены по названию встроенной функцией \code{sorted} (сравнение строк
        по кодам символов Unicode, поэтому, например, \code{İzmir} следует после \code{Zürich});
  \item каждое число выводится ровно с одним знаком после точки;
  \item среднее значение округляется до десятых по правилу <<половина вверх>>:
        $\lfloor 10\cdot\text{mean} + 0.5\rfloor / 10$, т.\,е. $1.25 \to 1.3$, $-1.25 \to -1.2$.
\end{itemize}

\subsubsection*{Материалы}

Каталог \code{data/} содержит:
\begin{itemize}
  \item \code{generate\_measurements.py}~— генератор входных файлов. Файлы объёмом в миллионы строк
        в репозиторий не добавляются и генерируются локально:
\begin{shell}
$ uv run data/generate_measurements.py 10_000_000 -o measurements_10m.txt
\end{shell}
        Генератор детерминирован: при одинаковых значениях \code{N} и \code{--seed} создаётся
        один и тот же файл. Ориентировочный размер: 10~млн строк~— около 140~МБ, 100~млн~— 1{,}4~ГБ,
        1~млрд~— 14~ГБ;
  \item \code{weather\_stations.csv}~— список станций со среднегодовыми температурами, используемый генератором;
  \item \code{measurements\_1k.txt}, \code{measurements\_100k.txt}~— небольшие входные файлы для отладки;
        \code{*.out}~— эталонные результаты для них;
  \item \code{check.py}~— программа проверки результата:
\begin{shell}
$ uv run main.py data/measurements_100k.txt \
    | python data/check.py data/measurements_100k.out
OK
\end{shell}
\end{itemize}

\subsubsection*{Требования к решению}

\begin{enumerate}
  \item \textbf{Программа} \code{main.py} принимает путь к файлу и выводит результат:
\begin{shell}
$ uv run main.py measurements_10m.txt
Abha;-0.8;18.3;48.0
Abidjan;-5.2;24.4;54.4
...
\end{shell}
  \item \textbf{Объём данных.} Программа обрабатывает файл из 100~млн строк (около 1{,}4~ГБ);
        объём используемой памяти не зависит от размера входного файла.
  \item \textbf{Ошибки.} При отсутствии аргументов или несуществующем файле программа выводит
        сообщение об ошибке в \code{stderr} и завершается с ненулевым кодом возврата, без вывода
        трассировки стека.
  \item \textbf{Корректность.} Результат совпадает с эталоном на файлах \code{measurements\_1k}
        и \code{measurements\_100k} (\code{check.py} выводит \code{OK}).
\end{enumerate}

%=====================================================================
\section{Порядок сдачи}

До защиты преподавателю направляется ссылка на репозиторий GitHub. Доступность ссылки
проверяется без входа в учётную запись (например, в приватном окне браузера); для приватного
репозитория проверяется наличие преподавателя в списке Collaborators. Сданное состояние
репозитория помечается тегом \code{lab1}; проверке подлежит именно оно:
\begin{shell}
$ git tag lab1
$ git push origin lab1
\end{shell}

На защите студент:
\begin{enumerate}
  \item демонстрирует репозиторий и историю коммитов; в репозитории отсутствуют \code{.venv}
        и файлы с данными большого объёма, присутствует \code{uv.lock};
  \item клонирует репозиторий в пустой каталог и запускает программу командой \code{uv run}
        без дополнительной настройки;
  \item запускает \code{check.py} на файле \code{measurements\_100k};
  \item отвечает на вопросы по коду.
\end{enumerate}

\subsection*{Возможные вопросы}

\begin{enumerate}
  \item Что такое коммит? Каково назначение команд \code{git add}, \code{git commit}, \code{git push}?
        Для чего нужен файл \code{.gitignore}?
  \item Для чего используются виртуальные окружения? Что содержит каталог \code{.venv} и почему
        он не добавляется в репозиторий? Чем \code{uv.lock} отличается от списка зависимостей
        в \code{pyproject.toml}?
  \item Каков результат выполнения \code{b = a; b.append(1)}, если \code{a}~— список?
        Каков результат \code{b = a; b += 1}, если \code{a}~— число? Чем объясняется различие?
  \item Чем различаются операторы \code{==} и \code{is}? В каких случаях уместно \code{is}?
  \item Какие объекты могут быть ключами словаря и почему?
  \item В чём состоит проблема изменяемого значения аргумента по умолчанию? Как её избежать?
  \item Проверяет ли интерпретатор аннотации типов во время выполнения? Для чего они используются?
  \item Чем различаются текстовый и бинарный режимы открытия файла? Каким образом в решении
        организовано чтение файла?
  \item Каким образом в решении обеспечено округление среднего значения по заданному правилу?
\end{enumerate}

\subsection*{Источники}
\begin{itemize}
  \item Официальное руководство по Python: \url{https://docs.python.org/3.13/tutorial/}
  \item Документация uv: \url{https://docs.astral.sh/uv/}
  \item PEP 8~— руководство по стилю кода: \url{https://peps.python.org/pep-0008/}
  \item Pro Git (книга о git, есть перевод на русский язык): \url{https://git-scm.com/book/ru/v2}
  \item One Billion Row Challenge: \url{https://github.com/gunnarmorling/1brc}
\end{itemize}

\end{document}
